\chapter{几何演化：输运恒等式与度量变分}

前七章中的几何结构都可以看成固定在一个流形上的对象。若对象随时间变化，必须先区分两种变化：一种只是由时间依赖微分同胚造成的“搬运”，另一种才是度量或其他几何数据本身的真实演化。把二者混在一起会产生大量看似不同的材料导数公式。本章从时间依赖拉回的精确微分公式出发，先统一所有输运方程，再对任意度量变分 \(\dot g:=\partial_tg\) 推出逆度量、体积、Levi--Civita 联络与 Hodge 星的变化；Ricci flow 只作为这些一般公式的一个特例。

\section{时间依赖微分同胚与精确输运公式}

设 \(\Phi_{t,s}:M\to M\) 是由时间依赖向量场 \(V_t\) 生成的两参数流：
\begin{equation}\label{eq:math-dg-time-dependent-flow}
\partial_t\Phi_{t,s}(p)
=V_t\bigl(\Phi_{t,s}(p)\bigr),
\qquad
\Phi_{s,s}=\IdMap.
\end{equation}
它满足组合律
\[
\Phi_{t,r}\circ\Phi_{r,s}=\Phi_{t,s}.
\]
若 \(T_t\) 是任意时间依赖张量场，则时间导数与拉回满足基本恒等式
\begin{equation}\label{eq:math-dg-transport-identity}
\boxed{
\frac{\dd}{\dd t}\bigl(\Phi_{t,s}^*T_t\bigr)
=\Phi_{t,s}^*
\bigl(\partial_tT_t+\Lie_{V_t}T_t\bigr).}
\end{equation}
这一定理是真正的几何版材料导数：\(\partial_t\) 比较同一坐标点上的显式时间变化，\(\Lie_{V_t}\) 补上由流动本身造成的搬运变化。

\begin{derivation}[时间依赖拉回的微分，即\eqnrefs{eq:math-dg-transport-identity}]
\emph{思路：}利用两参数流的组合律将 $\Phi_{t+\varepsilon,s}^*$ 分解为 $\Phi_{t,s}^*\circ\Phi_{t+\varepsilon,t}^*$，对短时流 $\Phi_{t+\varepsilon,t}^*$ 用 Lie 导数的一阶展开。

\emph{第一步：组合律分解。}固定 $t$，取小量 $\varepsilon$。两参数流的组合律给出
\begin{equation*}
 \Phi_{t+\varepsilon,s}
 =\Phi_{t+\varepsilon,t}\circ\Phi_{t,s}.
\end{equation*}
因此拉回分解为
\begin{equation*}
 \Phi_{t+\varepsilon,s}^*T_{t+\varepsilon}
 =\Phi_{t,s}^*
 \left(\Phi_{t+\varepsilon,t}^*T_{t+\varepsilon}\right).
\end{equation*}

\emph{第二步：短时流展开。}在时间 $t$ 附近，短时流 $\Phi_{t+\varepsilon,t}$ 由向量场 $V_t$ 生成，其拉回的一阶展开为
\begin{equation*}
 \Phi_{t+\varepsilon,t}^*
 =I+\varepsilon\,\Lie_{V_t}+o(\varepsilon),
\end{equation*}
而张量场的显式时间变化为
\begin{equation*}
 T_{t+\varepsilon}
 =T_t+\varepsilon\,\partial_tT_t+o(\varepsilon).
\end{equation*}

\emph{第三步：相乘并取差商。}
\begin{align*}
 \Phi_{t+\varepsilon,t}^*T_{t+\varepsilon}
 &=\bigl(I+\varepsilon\,\Lie_{V_t}+o(\varepsilon)\bigr)
  \bigl(T_t+\varepsilon\,\partial_tT_t+o(\varepsilon)\bigr)\\
 &=T_t+\varepsilon
  \bigl(\partial_tT_t+\Lie_{V_t}T_t\bigr)+o(\varepsilon).
\end{align*}
因此
\begin{align*}
 \frac{\dd}{\dd t}\bigl(\Phi_{t,s}^*T_t\bigr)
 &=\lim_{\varepsilon\to0}\frac{\Phi_{t+\varepsilon,s}^*T_{t+\varepsilon}-\Phi_{t,s}^*T_t}{\varepsilon}\\
 &=\Phi_{t,s}^*
 \bigl(\partial_tT_t+\Lie_{V_t}T_t\bigr),
\end{align*}
即\eqnrefs{eq:math-dg-transport-identity}。$\partial_t$ 比较同一坐标点上的显式时间变化，$\Lie_{V_t}$ 补上由流动本身造成的搬运变化——这是真正的几何版材料导数。
\end{derivation}

于是一般线性输运方程
\begin{equation}\label{eq:math-dg-geometric-transport-equation}
\partial_tT_t+\Lie_{V_t}T_t=S_t
\end{equation}
有形式精确解
\begin{equation}\label{eq:math-dg-geometric-transport-solution}
\boxed{
\Phi_{t,s}^*T_t
=T_s+\int_s^t\Phi_{\tau,s}^*S_\tau\,\dd\tau.}
\end{equation}
若 \(S_t=0\)，则 \(\Phi_{t,s}^*T_t=T_s\)。对微分形式还由 \(\dd\Phi^*=\Phi^*\dd\) 得到：闭性、恰当性与 de Rham 上同调类在纯输运中被保持。

\section{任意度量变分的代数部分}

现在固定底层流形，只令 Riemann 度量 \(g(t)\) 光滑依赖时间。记度量的时间导数为
\begin{equation}\label{eq:math-dg-metric-variation}
\dot g:=\partial_tg.
\end{equation}
它是一个对称二张量；任何具体几何流都只是给出 \(\dot g\) 的特殊表达。

逆度量的变化由恒等式 \(g^{ik}g_{kj}=\delta^i{}_j\) 决定：
\[
(\partial_tg^{ik})g_{kj}+g^{ik}\dot g_{kj}=0.
\]
乘以 \(g^{j\ell}\) 得
\begin{equation}\label{eq:math-dg-inverse-metric-variation}
\boxed{\partial_tg^{i\ell}=-\dot g^{i\ell}.}
\end{equation}
其中 \(\dot g^{i\ell}=g^{ik}g^{\ell j}\dot g_{kj}\) 在当前时间的度量下升指标。

体积形式的变化来自行列式。Jacobi 恒等式
\[
\partial_t\det g
=(\det g)\operatorname{tr}(g^{-1}\dot g)
\]
给出
\begin{equation}\label{eq:math-dg-volume-variation}
\boxed{
\partial_t\vol_g
=\frac12\operatorname{tr}_g(\dot g)\,\vol_g.}
\end{equation}
因此任何保持体积的一阶度量变分都必须满足 \(\operatorname{tr}_g \dot g=0\)；这比在某组坐标中直接对 \(\sqrt{\det g}\) 求导更具几何透明性。

\section{Levi--Civita 联络的变分}

每个 \(g(t)\) 都有唯一 Levi--Civita 联络 \(\nabla^{(t)}\)。固定时间无关的向量场 \(X,Y\)，定义
\[
\mathcal A(X,Y)
:=\partial_t\bigl(\nabla^{(t)}_XY\bigr).
\]
两个联络之差是张量，因此 \(\mathcal A\) 对 \(X,Y\) 都是函数线性的。它完全由 \(\dot g\) 的一阶协变导数确定。

\begin{derivation}[Levi--Civita 联络的一般变分公式]
在每个固定 \(t\) 上使用 Koszul 公式。向量场和 Lie 括号不随 \(t\) 变化，所以对时间求导时只有度量与联络贡献。对
\[
2g(\nabla_XY,Z)
=Xg(Y,Z)+Yg(Z,X)-Zg(X,Y)
+g([X,Y],Z)-g([Y,Z],X)+g([Z,X],Y)
\]
求导，左边为
\[
2\dot g(\nabla_XY,Z)+2g(\mathcal A(X,Y),Z).
\]
右边把每个 \(g\) 替换为 \(\dot g\)。再用当前 Levi--Civita 联络把普通方向导数拆成 \(\nabla \dot g\) 与对向量槽位的联络作用；所有含 \(\dot g(\nabla_XY,Z)\) 一类的项与左边对应项相消，得到
\begin{equation}\label{eq:math-dg-connection-variation}
\boxed{
\begin{aligned}
2g(\mathcal A(X,Y),Z)
={}&(\nabla_X\dot g)(Y,Z)
+(\nabla_Y\dot g)(X,Z)\\
&-(\nabla_Z\dot g)(X,Y).
\end{aligned}}
\end{equation}
坐标中即
\begin{equation}\label{eq:math-dg-christoffel-variation}
\boxed{
\partial_t\Gamma^k{}_{ij}
=\frac12g^{k\ell}
\left(
\nabla_i \dot g_{j\ell}
+\nabla_j \dot g_{i\ell}
-\nabla_\ell \dot g_{ij}
\right).}
\end{equation}

\emph{逐项消去的细节。}对右边第一项 $Xg(Y,Z)$ 求导，先用 $\dot g$ 替换 $g$ 得 $\dot g(Y,Z)$ 的方向导数 $X\dot g(Y,Z)$，再用 Levi--Civita 联络展开
\begin{equation*}
X\dot g(Y,Z)
=(\nabla_X\dot g)(Y,Z)
+\dot g(\nabla_XY,Z)+\dot g(Y,\nabla_XZ).
\end{equation*}
对其他两项作类似展开。含 $\dot g(\nabla_XY,Z)$ 的项与左边 $2\dot g(\nabla_XY,Z)$ 部分相消：右边三项给出 $+1+1+0=2$ 倍 $\dot g(\nabla_XY,Z)$（前两项含 $\nabla_XY$，第三项含 $\nabla_ZX$ 不贡献），与左边相消。剩余项整理即得协变导数形式。

\emph{张量性的直接验证。}要确认 $\mathcal A(X,Y)$ 是张量，需验证它对 $X,Y$ 都是函数线性的。对 $f\in C^\infty(M)$，
\begin{equation*}
\mathcal A(fX,Y)=\partial_t(\nabla_{fX}Y)=\partial_t(f\nabla_XY)=f\,\partial_t(\nabla_XY)=f\,\mathcal A(X,Y),
\end{equation*}
因 $f$ 与 $t$ 无关。同理对 $Y$。这印证了"两个联络之差是张量"的一般事实。

\emph{Ricci 流中的特例。}Ricci 流 $\partial_tg=-2\Ric$ 给出 $\dot g=-2\Ric$，代入得
\begin{equation*}
\partial_t\Gamma^k{}_{ij}
=-g^{k\ell}\left(
\nabla_i\Ric_{j\ell}
+\nabla_j\Ric_{i\ell}
-\nabla_\ell\Ric_{ij}
\right).
\end{equation*}
利用 contracted Bianchi 恒等式 $\nabla^\ell\Ric_{\ell j}=\tfrac12\nabla_j\mathcal R$ 与 deTurck 技巧，可把此式改写为
\begin{equation*}
\partial_t\Gamma^k{}_{ij}
=-g^{k\ell}\nabla_\ell\Ric_{ij}+\nabla_i(g^{k\ell}\Ric_{\ell j}),
\end{equation*}
第一项是 Ricci 张量的梯度，第二项可吸收为微分同胚作用。这就是 Ricci--deTurck 流的来源，它抛物化 Ricci 流，是证明短时间存在性的关键。

\emph{Hodge 星变分的预备。}Hodge 星同时依赖 $g^{-1}$ 与体积形式 $\vol_g$，两者都通过 $\dot g$ 变化。把 $\dot g$ 升指标得自伴 endomorphism $\mathcal V$，则
\begin{equation*}
\partial_t\vol_g=\tfrac12\operatorname{tr}\mathcal V\,\vol_g,
\qquad
\partial_tg^{-1}=-\mathcal V\circ g^{-1},
\end{equation*}
这两条与联络变分公式联合给出 Hodge 星的完整变分，是研究 harmonic form 在度量演化下行为的起点。
\end{derivation}

\eqnrefs{eq:math-dg-connection-variation}说明联络变化只含 \(\dot g\) 的一阶导数；因此曲率变化必然含 \(\dot g\) 的二阶导数。这一点可以直接写成精确的坐标无关公式。令 \(\mathcal A=\partial_t\nabla\)，则
\begin{equation}\label{eq:math-dg-curvature-variation-abstract}
\boxed{
\partial_tR(X,Y)Z
=(\nabla_X\mathcal A)(Y,Z)
-(\nabla_Y\mathcal A)(X,Z).}
\end{equation}
固定一组时间无关坐标，记 \(\mathcal A^\ell{}_{jk}=\partial_t\Gamma^\ell{}_{jk}\)，便有
\[
\partial_tR^\ell{}_{kij}
=\nabla_i\mathcal A^\ell{}_{jk}
-\nabla_j\mathcal A^\ell{}_{ik}.
\]
进一步对曲率算子取迹，得到
\begin{equation}\label{eq:math-dg-ricci-variation}
\partial_t\Ric_{jk}
=\nabla_i\mathcal A^i{}_{jk}
-\nabla_j\mathcal A^i{}_{ik}.
\end{equation}
把\eqnrefs{eq:math-dg-christoffel-variation}代入，并对 \(\mathcal R=g^{jk}\Ric_{jk}\) 求导，整理后得到标量曲率的一般变分
\begin{equation}\label{eq:math-dg-scalar-curvature-variation}
\boxed{
\partial_t\mathcal R
=-\langle\dot g,\Ric\rangle_g
+\nabla^i\nabla^j\dot g_{ij}
-\Delta_{\mathrm{LB}}(\operatorname{tr}_g\dot g).}
\end{equation}
这个公式仍未指定任何几何流；Ricci flow、Yamabe flow 或纯 Lie 导数变分都只是对 \(\dot g\) 的不同代入。

\begin{derivation}[由联络变分推出曲率与标量曲率变分]
从曲率定义
\[
R(X,Y)Z
=\nabla_X\nabla_YZ-\nabla_Y\nabla_XZ-\nabla_{[X,Y]}Z
\]
对时间求导。取 \(X,Y,Z\) 为时间无关向量场，且 \(\partial_t\nabla=\mathcal A\)，得到
\[
\begin{aligned}
\partial_tR(X,Y)Z
={}&\mathcal A(X,\nabla_YZ)+\nabla_X\mathcal A(Y,Z)\\
&-\mathcal A(Y,\nabla_XZ)-\nabla_Y\mathcal A(X,Z)
-\mathcal A([X,Y],Z).
\end{aligned}
\]
另一方面，
\[
(\nabla_X\mathcal A)(Y,Z)
=\nabla_X(\mathcal A(Y,Z))
-\mathcal A(\nabla_XY,Z)
-\mathcal A(Y,\nabla_XZ).
\]
对 \(X,Y\) 两项相减，并使用 Levi--Civita 联络无挠，
\[
\nabla_XY-\nabla_YX=[X,Y],
\]
便逐项恢复前一展开，从而得到\eqnrefs{eq:math-dg-curvature-variation-abstract}。在坐标基上取迹即得\eqnrefs{eq:math-dg-ricci-variation}。

最后对 \(\mathcal R=g^{jk}\Ric_{jk}\) 求导：
\[
\partial_t\mathcal R
=(\partial_tg^{jk})\Ric_{jk}
+g^{jk}\partial_t\Ric_{jk}.
\]
第一项由\eqnrefs{eq:math-dg-inverse-metric-variation}给出
\(-\langle\dot g,\Ric\rangle_g\)。第二项把\eqnrefs{eq:math-dg-christoffel-variation}代入\eqnrefs{eq:math-dg-ricci-variation}，利用 \(\nabla g=0\) 收缩指标，得到
\[
g^{jk}\partial_t\Ric_{jk}
=\nabla^i\nabla^j\dot g_{ij}
-\Delta_{\mathrm{LB}}(\operatorname{tr}_g\dot g),
\]
两项相加即为\eqnrefs{eq:math-dg-scalar-curvature-variation}。

\textbf{意义。}标量曲率变分公式 $\partial_t\mathcal R=\nabla^i\nabla^j\dot g_{ij}-\Delta_{\mathrm{LB}}(\operatorname{tr}_g\dot g)-\langle\dot g,\Ric\rangle_g$ 是几何流理论的基础。代入 Ricci flow $\dot g=-2\Ric$ 得 $\partial_t\mathcal R=\Delta\mathcal R+2|\Ric|^2$，由此推出标量曲率沿 Ricci flow 的极小值原理。代入 Einstein 方程的度规微扰 $\dot g=h$ 则给出广义相对论中曲率线性化——引力波方程的出发点。

\textbf{特例：共形变分。}取 $\dot g=2u\,g$（纯共形），则 $\operatorname{tr}_g\dot g=2nu$、$\nabla^i\nabla^j\dot g_{ij}=2\nabla^i\nabla_i u=2\Delta u$，代入得 $\partial_t\mathcal R=2\Delta u-2n\Delta u-2nu\mathcal R=-2(n-1)\Delta u-2nu\mathcal R$。在二维 ($n=2$) 化为 $\partial_t\mathcal R=-2\Delta u-4u\mathcal R$，即 Gauss 曲率在共形变分下的演化，是 Yamabe 问题的核心。
\end{derivation}


接下来考察总标量曲率泛函与 Einstein 张量。

一般曲率变分公式不仅服务于几何流，也直接给出 Riemann 几何中最重要的变分泛函之一。设 $M$ 紧致无边界，定义总标量曲率
\[
\mathcal S[g]
:=\int_M\mathcal R_g\,\vol_g.
\]
对任意对称二张量变分 $h=\dot g$，利用\eqnrefs{eq:math-dg-scalar-curvature-variation}和\eqnrefs{eq:math-dg-volume-variation}，有
\[
\begin{aligned}
\frac{\dd}{\dd t}\mathcal S[g(t)]
={}&\int_M
\left[
-\langle h,\Ric\rangle
+\nabla^i\nabla^jh_{ij}
-\Delta_{\mathrm{LB}}(\operatorname{tr}_gh)
\right]\vol_g\\
&+\frac12\int_M
\mathcal R\operatorname{tr}_g(h)\,\vol_g.
\end{aligned}
\]
闭流形上两个二阶导数项都是散度，积分为零，因此
\begin{equation}
\boxed{
\delta\mathcal S_g[h]
=-\int_M
\left\langle
\Ric-\frac12\mathcal R g,
h
\right\rangle_g\vol_g
}.
\label{eq:math-dg-einstein-hilbert-variation}
\end{equation}
括号中的
\[
G:=\Ric-\frac12\mathcal R g
\]
就是 Einstein tensor。它并不是人为组合出来的二阶张量，而是总标量曲率对度量的 $L^2$ 变分梯度（差一个符号）。

若只允许固定总体积的变分，需要加入约束
\[
\operatorname{Vol}(M,g)
=\int_M\vol_g=	ext{const.}
\]
并用 Lagrange multiplier。临界方程变成
\[
\Ric-\frac12\mathcal R g
+\lambda g=0.
\]
取迹可解出 $\lambda$ 与 $\mathcal R$ 的关系，最终等价于
\begin{equation}
\boxed{
\Ric=cg
}
\label{eq:math-dg-einstein-metric}
\end{equation}
其中 $c$ 为常数；这就是 Einstein metric。于是 Einstein 条件的变分来源先于任何特定坐标公式。

\begin{derivation}[固定体积约束为何给出 Einstein metric]
对体积泛函
\[
\mathcal V[g]=\int_M\vol_g
\]
有
\[
\delta\mathcal V_g[h]
=\frac12\int_M\operatorname{tr}_g(h)\,\vol_g
=\frac12\int_M\langle g,h\rangle_g\vol_g.
\]
临界条件
\[
\delta(\mathcal S+2\lambda\mathcal V)[h]=0
\quad\forall h
\]
给出
\[
-\Ric+\frac12\mathcal R g+\lambda g=0.
\]
在 $n$ 维中取迹：
\[
-\mathcal R+\frac n2\mathcal R+n\lambda=0,
\]
故 $\mathcal R$ 为常数。代回原方程，得到
\[
\Ric=\left(\frac12\mathcal R+\lambda\right)g=:cg.
\]
反过来任何 Einstein metric 满足带适当乘子的 Euler--Lagrange 方程。

\emph{体积变分的显式推导。}度量变分 $g_t=g+th$ 下体积形式的变化为
\begin{equation*}
\vol_{g_t}
=\sqrt{\frac{\det(g+th)}{\det g}}\,\vol_g
=\sqrt{\det(I+tg^{-1}h)}\,\vol_g.
\end{equation*}
对 $t$ 在 $t=0$ 求导，用 $\det(I+A)$ 的一阶展开 $\det(I+A)=1+\operatorname{tr}A+O(\|A\|^2)$，
\begin{equation*}
\left.\frac{\dd}{\dd t}\vol_{g_t}\right|_{t=0}
=\frac12\operatorname{tr}(g^{-1}h)\,\vol_g
=\frac12\operatorname{tr}_g(h)\,\vol_g.
\end{equation*}
积分即得体积泛函的变分公式。$\operatorname{tr}_g(h)=g^{ij}h_{ij}$ 是关于 $g$ 的迹，与坐标无关。

\emph{Einstein 常数的符号分类。}Einstein 方程 $\Ric=cg$ 中常数 $c$ 的符号给出三类典型几何：
\begin{itemize}
\item $c>0$：正曲率，如球面 $S^n$（$c=n-1$）与 Fano Kähler 流形；
\item $c=0$：Ricci 平坦，如平坦 $\RR^n$、Calabi--Yau 流形与 $G_2$ 流形，是弦论紧化的核心几何；
\item $c<0$：负曲率，如双曲空间 $H^n$（$c=-(n-1)$）与紧 Riemann 面的常负曲率度量。
\end{itemize}
Bishop--Gromov 体积比较定理给出 $c$ 对体积增长的精确控制：$c>0$ 给出多项式上界，$c=0$ 给出欧氏增长，$c<0$ 给出指数增长。

\emph{Kähler--Einstein 与 Calabi--Yau 定理。}在 Kähler 流形上，Einstein 条件 $\Ric=cg$ 等价于第一 Chern 类 $c_1(M)=c[\omega]$。当 $c_1=0$ 时 Calabi--Yau 定理保证 Ricci 平坦度量存在；当 $c_1>0$ 时存在性等价于 K-稳定性（Yau--Tian--Donaldson 猜想）；当 $c_1<0$ 时 Aubin--Yau 定理给出唯一 Kähler--Einstein 度量。三类情形的统一是几何不变量理论的核心问题。

\emph{Ricci 流的静态点。}Einstein 度量 $\Ric=cg$ 在 Ricci 流 $\partial_tg=-2\Ric$ 下演化为 $g(t)=(1-2ct)g$，是缩放固定点。因此 Einstein 度量是 Ricci 流的"理想终点"：Hamilton 的 Ricci 流程序把一般度量通过流推向 Einstein 度量，给出 Thurston 几何化猜想的解析证明路径。Perelman 的熵泛函单调性正是这一程序的关键工具。

\emph{物理建模：Einstein 真空方程。}广义相对论中 Einstein--Hilbert 作用量
\begin{equation*}
S_{\mathrm{EH}}[g]=\int_M\mathcal R\,\vol_g
\end{equation*}
的变分给出真空 Einstein 方程 $\Ric-\tfrac12\mathcal Rg=0$，取迹得 $\mathcal R=0$，故 $\Ric=0$。固定体积约束的变分问题对应宇宙常数 $\Lambda$ 的 Einstein 方程 $\Ric=\Lambda g$，描述 de Sitter 与 anti-de Sitter 时空。Einstein 度量因此既是 Riemann 几何的极值点，也是广义相对论的基本场方程解。
\end{derivation}

接下来考察Hodge 星的度量变分。

Hodge 星同时依赖逆度量和体积，因此它的变化不是只乘一个标量。令 \(\mathcal V:TM\to TM\) 为把 \(\dot g\) 升一个指标所得的自伴 endomorphism，即由
\begin{equation}\label{eq:math-dg-metric-variation-endomorphism}
g(\mathcal V X,Y)=\dot g(X,Y)
\end{equation}
定义的自伴 endomorphism。对 \(k\)-形式 \(\alpha\)，定义
\[
(\mathcal V^{[k]}\alpha)(X_1,\ldots,X_k)
:=\sum_{r=1}^k
\alpha(X_1,\ldots,\mathcal V X_r,\ldots,X_k).
\]
若 \(\alpha\) 本身不显含时间，则
\begin{equation}\label{eq:math-dg-hodge-star-variation}
\boxed{
\partial_t(*\alpha)
=*
\left[
\frac12\operatorname{tr}_g(\dot g)\,\alpha
-\mathcal V^{[k]}\alpha
\right].}
\end{equation}
若 \(\alpha=\alpha(t)\)，右端再加 \(*\partial_t\alpha\)。

\begin{derivation}[由 Hodge 星定义推出其变分，即\eqnrefs{eq:math-dg-hodge-star-variation}]
\emph{思路：}Hodge 星由度量通过内积和体积形式隐式定义。对度量变分，分别追踪内积变化和体积变化，再用 Hodge 星定义反读。

\emph{第一步：Hodge 星的隐式定义。}对任意时间无关 $k$-形式 $\alpha,\beta$，Hodge 星由
\begin{equation*}
 \alpha\wedge *\beta
 =\langle\alpha,\beta\rangle_g\,\vol_g
\end{equation*}
唯一确定。两边对 $t$ 求导，$\alpha$ 时间无关故只需变分 $*\beta$、$\langle\cdot,\cdot\rangle_g$ 和 $\vol_g$。

\emph{第二步：内积变分。}逆度量变化\eqnrefs{eq:math-dg-inverse-metric-variation}意味着 $k$-形式点态内积满足
\begin{equation*}
 \partial_t\langle\alpha,\beta\rangle_g
 =-\langle \mathcal V^{[k]}\alpha,\beta\rangle_g,
\end{equation*}
其中 $\mathcal V^{[k]}$ 是度量变分算子 $\mathcal V=g^{-1}\dot g$ 在 $k$-形式上的自然提升。

\emph{第三步：体积变分。}由\eqnrefs{eq:math-dg-volume-variation}，
\begin{equation*}
 \partial_t\vol_g
 =\frac{1}{2}\operatorname{tr}_g(\dot g)\,\vol_g.
\end{equation*}

\emph{第四步：合并。}对第一步等式求导，
\begin{align*}
 \alpha\wedge\partial_t(*\beta)
 &=\left[\partial_t\langle\alpha,\beta\rangle_g\right]\vol_g
  +\langle\alpha,\beta\rangle_g\,\partial_t\vol_g\\
 &=\left[
  -\langle \mathcal V^{[k]}\alpha,\beta\rangle_g
  +\frac{1}{2}\operatorname{tr}_g(\dot g)
  \langle\alpha,\beta\rangle_g
 \right]\vol_g.
\end{align*}

\emph{第五步：自伴性。}因为 $\mathcal V$ 自伴（$g^{-1}\dot g$ 关于 $g$ 对称），诱导的 $\mathcal V^{[k]}$ 在形式内积下也自伴：
\begin{equation*}
 \langle \mathcal V^{[k]}\alpha,\beta\rangle_g
 =\langle\alpha,\mathcal V^{[k]}\beta\rangle_g.
\end{equation*}

\emph{第六步：反读 Hodge 星。}将自伴性代入第四步结果，
\begin{equation*}
 \alpha\wedge\partial_t(*\beta)
 =\left\langle
  \frac{1}{2}\operatorname{tr}_g(\dot g)\,\alpha
  -\mathcal V^{[k]}\alpha,\,
  \beta
 \right\rangle_g\vol_g
 =\alpha\wedge *\left[
  \frac{1}{2}\operatorname{tr}_g(\dot g)\,\beta
  -\mathcal V^{[k]}\beta
 \right],
\end{equation*}
由 $\alpha$ 任意性即得到\eqnrefs{eq:math-dg-hodge-star-variation}。

\textbf{意义。}Hodge 星变分公式 $\dot{*}\alpha=*(\frac12\operatorname{tr}_g\dot g\,\alpha-\mathcal V^{[k]}\alpha)$ 度量了度规变化对 Hodge 对偶的影响。这在度规流理论中不可或缺：Ricci flow 下 Hodge 星随之演化，调和形式随之变化，影响上同调类的几何实现。实例：共形变分 $\dot g=2u\,g$ 下 $\operatorname{tr}_g\dot g=2nu$、$\mathcal V=0$，$\dot{*}\alpha=nu\,*\alpha$——Hodge 星按体积变化缩放，与共形因子 $e^{nu}$ 一致。
\end{derivation}

\section{纯微分同胚变化与 Ricci flow}

若度量变化只是由一族微分同胚产生，写
\[
g_t=\Psi_t^*g_0,
\qquad
W_t:=\partial_t\Psi_t\circ\Psi_t^{-1}.
\]
这里 \(W_t\) 是作用在目标流形上的 Eulerian 生成元。由拉回微分公式，
\[
\partial_tg_t
=\Psi_t^*(\Lie_{W_t}g_0).
\]
利用 Lie 导数对微分同胚的自然性，若把生成元搬回参考流形并记
\[
U_t:=\Psi_t^*W_t,
\]
则同一式也可写成
\begin{equation}\label{eq:math-dg-pure-diffeomorphism-metric-variation}
\boxed{\partial_tg_t=\Lie_{U_t}g_t.}
\end{equation}
由 Killing 方程推导中已经得到的恒等式，Levi--Civita 联络下
\begin{equation}\label{eq:math-dg-lie-metric-variation}
(\Lie_Ug)_{ij}
=\nabla_iU_j+\nabla_jU_i.
\end{equation}
这种 \(\dot g\) 只是在重新标记点，几何不变量随拉回一起变化，并不代表新的内在几何。这里必须区分 \(W_t\) 与 \(U_t=\Psi_t^*W_t\)；除非额外满足相应不变性，不能把 \(\Psi_t^*(\Lie_{W_t}g_0)\) 直接改写成 \(\Lie_{W_t}g_t\)。

真正的几何流则指定一个不纯粹来自 Lie 导数的 \(\dot g\)。最重要的例子之一是 Ricci flow
\begin{equation}\label{eq:math-dg-ricci-flow}
\partial_tg=-2\Ric(g).
\end{equation}
它只是一般变分公式中的特殊代入。由\eqnrefs{eq:math-dg-volume-variation}与标量曲率定义，立刻得到
\begin{equation}\label{eq:math-dg-ricci-flow-volume}
\partial_t\vol_g
=-\mathcal R\,\vol_g.
\end{equation}
由\eqnrefs{eq:math-dg-connection-variation}则得到 Levi--Civita 联络的变化由 \(\nabla\Ric\) 控制。将 \(\dot g=-2\Ric\) 代入\eqnrefs{eq:math-dg-scalar-curvature-variation}，先得到
\[
\partial_t\mathcal R
=2|\Ric|^2
-2\nabla^i\nabla^j\Ric_{ij}
+2\Delta_{\mathrm{LB}}\mathcal R.
\]
再用 contracted Bianchi 恒等式
\[
\nabla^i\Ric_{ij}=\frac12\nabla_j\mathcal R,
\]
得到
\begin{equation}\label{eq:math-dg-ricci-flow-scalar-curvature}
\boxed{
\partial_t\mathcal R
=\Delta_{\mathrm{LB}}\mathcal R+2|\Ric|^2.}
\end{equation}
因此 Ricci flow 下标量曲率的反应--扩散结构不是另行记忆的公式，而是一般度量变分的直接特例。

Ricci flow 仍具有微分同胚不变性，因此其坐标主符号带有 gauge 退化。选定固定背景度量 \(\bar g\)，定义 DeTurck 向量场
\begin{equation}\label{eq:math-dg-deturck-vector}
W^k
=g^{ij}\left(\Gamma^k{}_{ij}(g)-\bar\Gamma^k{}_{ij}(\bar g)\right).
\end{equation}
考虑 Ricci--DeTurck 方程
\begin{equation}\label{eq:math-dg-ricci-deturck-flow}
\boxed{
\partial_tg=-2\Ric(g)+\Lie_Wg.}
\end{equation}
在局部坐标中，\(-2\Ric\) 的非严格抛物部分恰被 \(\Lie_Wg\) 的最高阶项抵消，留下
\[
\partial_tg_{ij}
=g^{ab}\partial_a\partial_bg_{ij}
+\text{至多一阶导数的非线性项}.
\]
所以固定 \(\bar g\) 后，Ricci--DeTurck 系统具有严格抛物主部。求解它以后，再沿由 \(-W\) 生成的时间依赖微分同胚拉回，就恢复 Ricci flow。DeTurck trick 因而只是固定微分同胚 gauge，并没有改变内在几何内容。


Ricci flow 会改变总体积。若希望比较固定体积下的形状演化，可定义平均标量曲率
\[
\bar{\mathcal R}(t)
:=\frac{1}{\operatorname{Vol}(M,g(t))}
\int_M\mathcal R\,\vol_g
\]
并考虑 normalized Ricci flow
\begin{equation}
\boxed{
\partial_tg
=-2\Ric+\frac{2}{n}\bar{\mathcal R}\,g
}.
\label{eq:math-dg-normalized-ricci-flow}
\end{equation}
由体积变分公式，
\[
\frac{\dd}{\dd t}\operatorname{Vol}(M,g)
=\int_M
\left(-\mathcal R+\bar{\mathcal R}\right)\vol_g=0.
\]
因此归一化项不是随意加入的尺度修正，而是精确消去全局体积漂移的 Lagrange-multiplier 型项。

对标量曲率，代入一般变分公式得到
\begin{equation}
\partial_t\mathcal R
=\Delta_{\mathrm{LB}}\mathcal R
+2|\Ric|^2
-\frac{2}{n}\bar{\mathcal R}\,\mathcal R.
\label{eq:math-dg-normalized-ricci-scalar}
\end{equation}
若在某个时间区间内曲率有下界或上界，这类反应--扩散方程可以直接配合 maximum principle。对非归一化 Ricci flow，若 $\mathcal R$ 在时刻 $t$ 取得空间最小值，则在该点 $\Delta_{\mathrm{LB}}\mathcal R\ge0$，并由
\[
|\Ric|^2\ge\frac1n\mathcal R^2
\]
得到
\[
\partial_t\mathcal R_{\min}
\ge\frac{2}{n}\mathcal R_{\min}^2
\]
的 barrier 意义估计。它说明 scalar curvature 的负下界不会任意快向 $-\infty$ 漂移，而正曲率又可能在有限时间内因 reaction term 增强。这类 maximum-principle 推导是几何流从张量 PDE 获得几何结论的基本方式。


接下来考察Ricci tensor 演化、尺度律与奇性判据。

标量曲率只记录曲率张量的一次收缩。若要研究 Ricci flow 是否形成奇性，必须继续跟踪 $\Ric$ 和完整 $\operatorname{Rm}$。把 $h=-2\Ric$ 代入一般 Ricci 变分式，并利用两次 contracted Bianchi 恒等式整理，可得 $(0,2)$ 型 Ricci tensor 的精确演化
\begin{equation}
\boxed{
\partial_t\Ric_{ij}
=\Delta_{\mathrm{LB}}\Ric_{ij}
+2R_{ikj\ell}\Ric^{k\ell}
-2\Ric_i{}^k\Ric_{kj}.
}
\label{eq:math-dg-ricci-flow-ricci-evolution}
\end{equation}
对它取迹时要同时考虑逆度量的变化
\[
\partial_tg^{ij}=2\Ric^{ij},
\]
于是恰好恢复\eqnrefs{eq:math-dg-ricci-flow-scalar-curvature}。因此标量曲率演化并不是独立公式，而是 Ricci tensor 演化与“收缩本身也随时间变化”共同产生的结果。

完整 Riemann tensor 的分量公式较长，但其结构非常稳定：
\begin{equation}
\boxed{
(\partial_t-\Delta_{\mathrm{LB}})\operatorname{Rm}
=\operatorname{Rm}*\operatorname{Rm},
}
\label{eq:math-dg-ricci-flow-rm-schematic}
\end{equation}
其中 $*$ 表示由度量收缩得到的有限个二次曲率组合。相应地，曲率范数满足
\begin{equation}
(\partial_t-\Delta_{\mathrm{LB}})|\operatorname{Rm}|^2
=-2|\nabla\operatorname{Rm}|^2+\operatorname{Rm}*\operatorname{Rm}*\operatorname{Rm}.
\label{eq:math-dg-ricci-flow-rm-norm}
\end{equation}
真正重要的是右边只有“耗散梯度项 + 三次反应项”。这与半线性抛物 PDE 的 blow-up 结构完全一致：扩散试图平滑曲率，非线性曲率自相互作用则可能放大高曲率区。

Ricci flow 具有精确抛物尺度。若 $g(t)$ 是解，则对任意 $\lambda>0$，
\begin{equation}
\widetilde g(t)
:=\lambda g(t/\lambda)
\label{eq:math-dg-ricci-parabolic-rescaling}
\end{equation}
仍是 Ricci flow。长度尺度变为 $\sqrt\lambda$，时间尺度变为 $\lambda$，而
\[
|\widetilde{\operatorname{Rm}}|_{\widetilde g}
=\lambda^{-1}|\operatorname{Rm}|_g.
\]
所以在高曲率点附近选择与曲率大小相匹配的抛物缩放，就能把该点曲率归一到 $O(1)$。几何流奇性分析中的 tangent flow、ancient solution 与 blow-up limit 都来自这一尺度律，而不是额外的技巧。

对闭流形，Ricci--DeTurck 的严格抛物性给出短时存在唯一性。更进一步，若最大存在区间为 $[0,T)$ 且 $T<\infty$，则 Hamilton extension criterion 说明
\begin{equation}
\boxed{
\limsup_{t\uparrow T}
\sup_{x\in M}|\operatorname{Rm}(x,t)|=\infty.
}
\label{eq:math-dg-ricci-extension-criterion}
\end{equation}
换言之，有限时间失效不能只由坐标或微分同胚 gauge 造成；必定有真正的内在曲率 blow-up。这个结论与非线性 PDE 章节中的 continuation criterion 具有相同逻辑：只要控制决定局部适定性的最高阶几何量，就能把解延拓过当前时间。

二维情形把这些结构压缩得最透明。若 $n=2$，则
\[
\Ric=\frac12\mathcal R g,
\]
故 Ricci flow 变成
\begin{equation}
\partial_tg=-\mathcal R g,
\label{eq:math-dg-ricci-flow-two-dimensional-metric}
\end{equation}
而标量曲率满足
\begin{equation}
\boxed{
\partial_t\mathcal R
=\Delta_{\mathrm{LB}}\mathcal R+\mathcal R^2.
}
\label{eq:math-dg-ricci-flow-two-dimensional-scalar}
\end{equation}
这说明二维 Ricci flow 本质上是一个带几何耦合的反应--扩散方程；高维中新出现的困难正来自 Weyl curvature 与 Ricci tensor 不再由一个标量决定。

接下来考察Ricci soliton、自相似解与熵型泛函。

Ricci flow 的 stationary point 不能只写成 $\Ric=0$，因为几何流允许同时发生整体缩放与微分同胚重标记。把这两个“平凡运动”都除掉后，自相似解由 Ricci soliton 方程刻画：存在向量场 $X$ 与常数 $\lambda$ 使
\begin{equation}
\boxed{
\Ric+\frac12\Lie_Xg=\lambda g.
}
\label{eq:math-dg-ricci-soliton}
\end{equation}
$\lambda>0,=0,<0$ 时分别称 shrinking、steady、expanding soliton。若
\[
X=\nabla f,
\]
则
\[
\frac12\Lie_{\nabla f}g=\nabla^2f,
\]
于是得到 gradient Ricci soliton
\begin{equation}
\boxed{
\Ric+\nabla^2f=\lambda g.
}
\label{eq:math-dg-gradient-ricci-soliton}
\end{equation}

\begin{derivation}[soliton 方程为何精确产生自相似 Ricci flow]
设 $(g,X,\lambda)$ 满足\eqnrefs{eq:math-dg-ricci-soliton}，令
\[
c(t)=1-2\lambda t.
\]
在 $c(t)>0$ 的区间上，令 $\phi_t$ 由时间依赖向量场
\[
Y_t=\frac{1}{c(t)}X
\]
生成，并定义
\begin{equation}
g(t)=c(t)\phi_t^*g.
\label{eq:math-dg-soliton-selfsimilar-flow}
\end{equation}
则
\begin{align*}
\partial_tg(t)
&=c'(t)\phi_t^*g
+c(t)\phi_t^*(\Lie_{Y_t}g)\\
&=\phi_t^*(-2\lambda g+\Lie_Xg)\\
&=-2\phi_t^*\Ric_g.
\end{align*}
常数缩放不改变 $(0,2)$ 型 Ricci tensor，而 Ricci tensor 对微分同胚自然，因此
\[
\Ric_{g(t)}=\phi_t^*\Ric_g.
\]
所以
\[
\partial_tg(t)=-2\Ric_{g(t)}.
\]
这证明 soliton 不是"近似自相似"，而是 Ricci flow 在缩放与微分同胚作用下的精确固定点。

\emph{三类 soliton 的分类。}根据 $\lambda$ 的符号，soliton 分为三类：
\begin{itemize}
\item \emph{Shrinker}（$\lambda>0$）：流在 $t\to1/(2\lambda)$ 时收缩到一点，是 Ricci 流奇点的 blow-up 模型。典型例子是 Gauss shrinker $g_{\mathrm{euc}}$、$X=-\tfrac12\nabla|x|^2$，对应 Gaussian 稳定分布。
\item \emph{Steady}（$\lambda=0$）：流仅由微分同胚驱动，体积不变。典型例子是 cigar soliton $g=\frac{\dd x^2+\dd y^2}{1+x^2+y^2}$，对应稳态热核。
\item \emph{Expander}（$\lambda<0$）：流对所有 $t<0$ 存在，从 $t\to-\infty$ 的一点展开。典型例子是 expanding sphere，对应热核扩散。
\end{itemize}

\emph{Gradient soliton 的简化。}当 $X=\nabla f$ 为梯度场时，soliton 方程简化为
\begin{equation*}
\Ric+\nabla^2f=\lambda g.
\end{equation*}
取迹得 $\mathcal R+\Delta f=n\lambda$。对\eqnrefs{eq:math-dg-gradient-ricci-soliton}取散度，用 contracted Bianchi 恒等式 $\nabla^i\Ric_{ij}=\tfrac12\nabla_j\mathcal R$ 与 Hessian 散度公式
\begin{equation*}
\operatorname{div}(\nabla^2f)_j=\nabla_j\Delta f+\Ric_{ij}\nabla^if,
\end{equation*}
得到
\begin{equation*}
\nabla_j\mathcal R+\nabla_j\Delta f+\Ric_{ij}\nabla^if=n\lambda\nabla_jf,
\end{equation*}
结合迹方程给出 $\nabla\mathcal R=2\Ric(\nabla f,\cdot)$，积分得 first integral
\begin{equation*}
\mathcal R+|\nabla f|^2-2\lambda f=C.
\end{equation*}
这是 Perelman $\mathcal W$ 泛函的临界点方程。

\emph{Perelman 熵与 soliton。}Perelman $\mathcal W$ 泛函
\begin{equation*}
\mathcal W[g,f,\tau]
=\int_M\left[\tau(\mathcal R+|\nabla f|^2)-f-n\right]\frac{\ee^{-f}}{(4\pi\tau)^{n/2}}\vol_g
\end{equation*}
在 Ricci 流耦合 $\partial_tf=-\Delta f-\mathcal R+\frac n{2\tau}$、$\partial_t\tau=-1$ 下单调不增，临界点恰为 gradient shrinker。因此 shrinker 是 Ricci 流的"熵极大点"，对应统计力学中的高斯分布。

\emph{举例：Gauss shrinker 与雪茄 soliton。}在 $\RR^n$ 上取 $g_{\mathrm{euc}}$、$f=\tfrac{|x|^2}{4\tau}$、$\lambda=\tfrac1{2\tau}$，直接验证
\begin{equation*}
\Ric=0,\qquad
\nabla^2f=\frac{1}{2\tau}g_{\mathrm{euc}},
\end{equation*}
满足 shrinker 方程。这给出 Ricci 流从一点展开的标准 Euclidean 解，对应热方程基本解。二维雪茄 soliton $g=\frac{\dd x^2+\dd y^2}{1+x^2+y^2}$、$X=-2(x\partial_x+y\partial_y)$、$\lambda=0$ 是 steady soliton，曲率 $\mathcal R=\frac{4}{1+r^2}$ 在无穷远趋于零，是二维 Ricci 流分类的基本模型。

\emph{物理建模：黑洞几何与 soliton。}Bryant soliton 是三维 Ricci 流中非紧 shrinker 的典型例子，其渐近行为是圆柱形，对应正曲率 Ricci 流的奇点模型。在广义相对论中，gradient Ricci soliton 方程与静态 Einstein 方程通过 conformal 变换相关，给出某些静态时空的几何描述，例如 Schwarzschild 几何在适当坐标下满足 soliton 型方程。
\end{derivation}

对 gradient soliton，取迹得到
\begin{equation}
\mathcal R+\Delta_{\mathrm{LB}}f=n\lambda.
\label{eq:math-dg-gradient-soliton-trace}
\end{equation}
对\eqnrefs{eq:math-dg-gradient-ricci-soliton} 取散度，并使用 contracted Bianchi 恒等式与 Hessian 的交换公式，可得
\[
\dd\mathcal R=2\Ric(\nabla f,\cdot).
\]
于是
\begin{equation}
\boxed{
\mathcal R+|\nabla f|^2-2\lambda f=C
}
\label{eq:math-dg-gradient-soliton-first-integral}
\end{equation}
在每个连通分支上为常数。这个 first integral 在研究 shrinker 的加权体积、势函数增长与 blow-up limit 时非常重要。

自相似结构还可以由一个单调泛函识别。对闭流形定义 Perelman $\mathcal F$-functional
\begin{equation}
\mathcal F(g,f)
:=\int_M
\left(\mathcal R+|\nabla f|^2\right)
\ee^{-f}\vol_g,
\qquad
\int_M\ee^{-f}\vol_g=1.
\label{eq:math-dg-perelman-f-functional}
\end{equation}
若 $(g,f)$ 满足耦合系统
\begin{equation}
\partial_tg=-2\Ric,
\qquad
\partial_tf
=-\Delta_{\mathrm{LB}}f+|\nabla f|^2-\mathcal R,
\label{eq:math-dg-perelman-coupled-flow}
\end{equation}
则加权测度总质量保持，并且
\begin{equation}
\boxed{
\frac{\dd}{\dd t}\mathcal F(g,f)
=2\int_M
|\Ric+\nabla^2f|^2
\ee^{-f}\vol_g
\ge0.
}
\label{eq:math-dg-perelman-f-monotonicity}
\end{equation}
因此 Ricci flow 并不是一个没有 Lyapunov 结构的任意非线性 PDE；在与适当 backward heat-type 势函数耦合后，它具有严格的熵型单调量。

先验证约束确实保持。由\eqnrefs{eq:math-dg-ricci-flow-volume} 与\eqnrefs{eq:math-dg-perelman-coupled-flow}，
\begin{align*}
\frac{\dd}{\dd t}\int_M\ee^{-f}\vol_g
&=\int_M(-\partial_tf-\mathcal R)\ee^{-f}\vol_g\\
&=\int_M
(\Delta_{\mathrm{LB}}f-|\nabla f|^2)
\ee^{-f}\vol_g=0,
\end{align*}
其中最后一步用分部积分
\[
\int_M\ee^{-f}\Delta_{\mathrm{LB}}f\,\vol_g
=\int_M\ee^{-f}|\nabla f|^2\vol_g.
\]
而\eqnrefs{eq:math-dg-perelman-f-monotonicity} 则来自标量曲率、梯度能与加权体积三项变分的精确抵消；剩余项完成平方为
\[
2|\Ric+\nabla^2f|^2.
\]
这与前面“总标量曲率的一阶变分给 Einstein tensor”形成对应：这里不是寻找静态 Einstein 临界点，而是在加权几何中找到 Ricci flow 的梯度型结构。

对 shrinker 的完整刻画还需要加入尺度归一化后的 $\mathcal W$-entropy；本章先保留到 $\mathcal F$ 这一层，因为它已足以说明三个核心事实：diffeomorphism gauge 由 DeTurck 处理，自相似 blow-up 模型由 Ricci soliton 描述，而熵型单调量为排除某些回返行为和分析奇性极限提供全局方向。

\begin{exbox}[闭形式的冻结输运]
设 \(\alpha_t\in\Omega^k(M)\) 满足
\[
\partial_t\alpha_t+\Lie_{V_t}\alpha_t=0.
\]
由\eqnrefs{eq:math-dg-geometric-transport-solution}，
\[
\Phi_{t,s}^*\alpha_t=\alpha_s.
\]
若 \(\dd\alpha_s=0\)，则
\[
\Phi_{t,s}^*(\dd\alpha_t)
=\dd(\Phi_{t,s}^*\alpha_t)
=\dd\alpha_s=0.
\]
微分同胚拉回可逆，所以 \(\dd\alpha_t=0\)。这就是理想流体和磁流体中“冻结形式”结论的几何母结构。
\end{exbox}
